% GENERATED FILE. Edit canonical lessons, then run npm run generate:books.
% canonical-source-hash: fnv1a64:b2a1a3b3fcff5dd9
% canonical-lessons: JA-C493-reitouko, JA-C493-todana, JA-C493-suihanki, JA-C493-araimono, JA-C493-reitou-suru

\chapter{Freezer and Cupboard}
\label{ch:ja-freezer-and-cupboard}
\hlchaptermodality{493}

\begin{chapteropening}
\textbf{By the end of this chapter:} \emph{I can say where food and dishes are kept.}

Complete the last lesson of chapter 493: I can say where food and dishes are kept.
\end{chapteropening}

\section[\texorpdfstring{\ja{れいとうこ} (reitōko)}{れいとうこ (reitōko)}]{\ja{れいとうこ} (reitōko) — a freezer}

\label{lesson:JA-C493-reitouko}



\begin{quote}
Before the new one: say the Japanese for to steam (food), then the Japanese for to peel.
\end{quote}

\subsection*{You'll want to know: \ja{れいとうこ}}
\textbf{\ja{れいとうこ}} — \emph{reitōko} — \textquotedblleft{}a freezer\textquotedblright{}.

\begin{quote}
\textbf{\ja{れいとうこに} \ja{こおりが} \ja{あります。}} — \emph{reitōko ni kōri ga arimasu} — \textquotedblleft{}There is ice in the freezer.\textquotedblright{}
\end{quote}

\begin{cousinweb}
\textbf{\ja{れいとうこ}} (\emph{reitōko}) is Sino-Japanese (kango): three Chinese characters read with their on'yomi, readings borrowed from Chinese. They are \emph{rei}, \textquotedblleft{}cool\textquotedblright{}, \emph{tō}, \textquotedblleft{}frozen\textquotedblright{} (Middle Chinese \emph{tuwngH}) and \emph{ko}, \textquotedblleft{}warehouse\textquotedblright{} (\emph{khuH}) (KANJIDIC2; Middle Chinese in Baxter's transcription).
\end{cousinweb}

\begin{grammarlens}[title={What you've built}]
One more word for describing what you lived and felt.
\end{grammarlens}

\subsection*{Your turn}
\begin{itemize}
\raggedright
  \item \emph{Say it:} \emph{reitōko}
  \item \emph{Say it:} \emph{reitōko}, once more
  \item \emph{Say it:} \emph{reitōko ni kōri ga arimasu} — There is ice in the freezer.
  \item \emph{Recall:} say the Japanese for to steam (food), then the Japanese for to peel, then say \emph{reitōko} again
\end{itemize}

\begin{tcolorbox}[breakable,colback=teal!4,colframe=teal!35!black,title={Before you move on}]
What does \ja{れいとうこ} mean? (\textquotedblleft{}A freezer\textquotedblright{}.) Say it once more.
\end{tcolorbox}
\section[\texorpdfstring{\ja{とだな} (todana)}{とだな (todana)}]{\ja{とだな} (todana) — a cupboard}

\label{lesson:JA-C493-todana}



\begin{quote}
Before the new one: say the Japanese for to peel, then the Japanese for a freezer.
\end{quote}

\subsection*{You'll want to know: \ja{とだな}}
\textbf{\ja{とだな}} — \emph{todana} — \textquotedblleft{}a cupboard\textquotedblright{}.

\begin{quote}
\textbf{\ja{ちゃわんを} \ja{とだなに} \ja{しまいます。}} — \emph{chawan o todana ni shimaimasu} — \textquotedblleft{}I put the teacups away in the cupboard.\textquotedblright{}
\end{quote}

\begin{cousinweb}
\textbf{\ja{とだな}} (\emph{todana}) is a native Japanese compound (wago), inherited rather than borrowed from Chinese: built from \emph{to}, \textquotedblleft{}door\textquotedblright{} and \emph{dana}, \textquotedblleft{}shelf\textquotedblright{} (JMdict, KANJIDIC2).
\end{cousinweb}

\begin{grammarlens}[title={What you've built}]
One more word for describing what you lived and felt.
\end{grammarlens}

\subsection*{Your turn}
\begin{itemize}
\raggedright
  \item \emph{Say it:} \emph{todana}
  \item \emph{Say it:} \emph{todana}, once more
  \item \emph{Say it:} \emph{chawan o todana ni shimaimasu} — I put the teacups away in the cupboard.
  \item \emph{Recall:} say the Japanese for to peel, then the Japanese for a freezer, then say \emph{todana} again
\end{itemize}

\begin{tcolorbox}[breakable,colback=teal!4,colframe=teal!35!black,title={Before you move on}]
What does \ja{とだな} mean? (\textquotedblleft{}A cupboard\textquotedblright{}.) Say it once more.
\end{tcolorbox}
\section[\texorpdfstring{\ja{すいはんき} (suihanki)}{すいはんき (suihanki)}]{\ja{すいはんき} (suihanki) — a rice cooker}

\label{lesson:JA-C493-suihanki}



\begin{quote}
Before the new one: say the Japanese for a freezer, then the Japanese for a cupboard.
\end{quote}

\subsection*{You'll want to know: \ja{すいはんき}}
\textbf{\ja{すいはんき}} — \emph{suihanki} — \textquotedblleft{}a rice cooker\textquotedblright{}.

\begin{quote}
\textbf{\ja{すいはんきで} \ja{ごはんを} \ja{たきます。}} — \emph{suihanki de gohan o takimasu} — \textquotedblleft{}I cook rice in the rice cooker.\textquotedblright{}
\end{quote}

\begin{cousinweb}
\textbf{\ja{すいはんき}} (\emph{suihanki}) is Sino-Japanese (kango): three Chinese characters read with their on'yomi, readings borrowed from Chinese. They are \emph{sui}, \textquotedblleft{}cook\textquotedblright{} (Middle Chinese \emph{tsyhwe}), \emph{han}, \textquotedblleft{}meal\textquotedblright{} (\emph{bjonH}) and \emph{ki}, \textquotedblleft{}utensil\textquotedblright{} (\emph{khijH}) (KANJIDIC2; Middle Chinese in Baxter's transcription).
\end{cousinweb}

\begin{grammarlens}[title={What you've built}]
One more word for describing what you lived and felt.
\end{grammarlens}

\subsection*{Your turn}
\begin{itemize}
\raggedright
  \item \emph{Say it:} \emph{suihanki}
  \item \emph{Say it:} \emph{suihanki}, once more
  \item \emph{Say it:} \emph{suihanki de gohan o takimasu} — I cook rice in the rice cooker.
  \item \emph{Recall:} say the Japanese for a freezer, then the Japanese for a cupboard, then say \emph{suihanki} again
\end{itemize}

\begin{tcolorbox}[breakable,colback=teal!4,colframe=teal!35!black,title={Before you move on}]
What does \ja{すいはんき} mean? (\textquotedblleft{}A rice cooker\textquotedblright{}.) Say it once more.
\end{tcolorbox}
\section[\texorpdfstring{\ja{あらいもの} (araimono)}{あらいもの (araimono)}]{\ja{あらいもの} (araimono) — the washing-up; things to be washed}

\label{lesson:JA-C493-araimono}



\begin{quote}
Before the new one: say the Japanese for a cupboard, then the Japanese for a rice cooker.
\end{quote}

\subsection*{You'll want to know: \ja{あらいもの}}
\textbf{\ja{あらいもの}} — \emph{araimono} — \textquotedblleft{}the washing-up; things to be washed\textquotedblright{}.

\begin{quote}
\textbf{\ja{ばんごはんの} \ja{あとで} \ja{あらいものを} \ja{します。}} — \emph{bangohan no ato de araimono o shimasu} — \textquotedblleft{}After dinner I do the washing-up.\textquotedblright{}
\end{quote}

\begin{cousinweb}
\textbf{\ja{あらいもの}} (\emph{araimono}) is a native Japanese compound (wago), inherited rather than borrowed from Chinese: built from \emph{arai}, \textquotedblleft{}washing\textquotedblright{}, from \emph{arau}, \textquotedblleft{}to wash\textquotedblright{}, and \emph{mono}, \textquotedblleft{}thing\textquotedblright{} (JMdict, KANJIDIC2).
\end{cousinweb}

\begin{grammarlens}[title={What you've built}]
One more word for describing what you lived and felt.
\end{grammarlens}

\subsection*{Your turn}
\begin{itemize}
\raggedright
  \item \emph{Say it:} \emph{araimono}
  \item \emph{Say it:} \emph{araimono}, once more
  \item \emph{Say it:} \emph{bangohan no ato de araimono o shimasu} — After dinner I do the washing-up.
  \item \emph{Recall:} say the Japanese for a cupboard, then the Japanese for a rice cooker, then say \emph{araimono} again
\end{itemize}

\begin{tcolorbox}[breakable,colback=teal!4,colframe=teal!35!black,title={Before you move on}]
What does \ja{あらいもの} mean? (\textquotedblleft{}The washing-up; things to be washed\textquotedblright{}.) Say it once more.
\end{tcolorbox}
\section[\texorpdfstring{\ja{れいとうする} (reitō suru)}{れいとうする (reitō suru)}]{\ja{れいとうする} (reitō suru) — to freeze (food)}

\label{lesson:JA-C493-reitou-suru}



\begin{quote}
Before the new one: say the Japanese for a rice cooker, then the Japanese for the washing-up.
\end{quote}

\subsection*{You'll want to know: \ja{れいとうする}}
\textbf{\ja{れいとうする}} — \emph{reitō suru} — \textquotedblleft{}to freeze (food)\textquotedblright{}.

\begin{quote}
\textbf{\ja{にくが} \ja{のこったら、れいとうします。}} — \emph{niku ga nokottara, reitōshimasu} — \textquotedblleft{}If there is meat left over, I freeze it.\textquotedblright{}
\end{quote}

Polite \textbf{\ja{れいとうします}}, \emph{reitōshimasu}; a \ja{する} verb. Past (polite): \textbf{\ja{ははが} \ja{さかなを} \ja{れいとうしました。}}, \emph{haha ga sakana o reitōshimashita}.

\begin{cousinweb}
\textbf{\ja{れいとうする}} (\emph{reitō suru}) is the Sino-Japanese (kango) noun \emph{reitō} with the native verb \emph{suru}, \textquotedblleft{}to do\textquotedblright{}: two Chinese characters read with their on'yomi, readings borrowed from Chinese. They are \emph{rei}, \textquotedblleft{}cool\textquotedblright{} and \emph{tō}, \textquotedblleft{}frozen\textquotedblright{} (Middle Chinese \emph{tuwngH}) (KANJIDIC2; Middle Chinese in Baxter's transcription).
\end{cousinweb}

\begin{grammarlens}[title={What you've built}]
One more word for describing what you lived and felt.
\end{grammarlens}

\subsection*{Your turn}
\begin{itemize}
\raggedright
  \item \emph{Say it:} \emph{reitō suru}
  \item \emph{Say it:} \emph{reitō suru}, once more
  \item \emph{Say it:} \emph{niku ga nokottara, reitōshimasu} — If there is meat left over, I freeze it.
  \item \emph{Recall:} say the Japanese for a rice cooker, then the Japanese for the washing-up, then say \emph{reitō suru} again
\end{itemize}

\begin{tcolorbox}[breakable,colback=teal!4,colframe=teal!35!black,title={Before you move on}]
What does \ja{れいとうする} mean? (\textquotedblleft{}To freeze (food)\textquotedblright{}.) Say it once more.
\end{tcolorbox}
